\ExampleChapter{摘\quad 要}{摘要}
无人机多变量飞行记录为时序建模与异常分析提供了典型的数据组织场景。
为演示学位论文从问题定义、算法描述到实验展示的完整排版过程，
本文构建了一个基于合成飞行数据的预测残差检测案例。
该案例不依赖真实无人机平台，也不声称提出新的检测算法，
主要用于验证 GZUThesis 对复杂学术文档元素的支持能力。

首先，利用具有周期项、变量耦合项与随机扰动的函数生成六维时序数据，
并在测试架次中注入渐变漂移、高频振动和偏置三类异常。
其次，以完整架次为单位划分训练集、校准集与测试集，
仅使用训练集计算标准化统计量和拟合岭回归预测器。
随后，将一步预测误差聚合为异常分数，采用指数加权移动平均抑制局部波动，
并使用正常校准数据的分位数确定阈值。
最后，与持续值预测、滑动平均预测及未平滑的岭回归分数进行比较。

在固定随机种子的本次合成实验中，岭回归与平滑组合的逐点精确率、召回率和
F1 值分别为 \DemoPrecision\%、\DemoRecall\% 和 \DemoFone\%。
这些数值只描述随附脚本生成的数据，不代表真实飞行故障的检测能力。
实例同时说明，平滑并不保证检测性能提升；模型、阈值和评价规则应作为独立因素记录。
通过将图表、统计量及参考文献纳入可复现工程，
可以减少论文编辑过程中编号、引用和实验文字不一致的问题。

\par\vspace{1em}
\noindent\textbf{关键词：}无人机；时间序列；异常检测；预测残差；可复现排版
